% Nová příprava: 2026-10-28; 1. ročník SPŠ, etapa 13–24.
% Rozpis: 1_rocnik/tematicke_plany/Rozpis_vyuky_MAT_1_rocnik.tex, hodina 15.
% Původ: vlastní formulace výkladu a vlastní úlohy; metodická inspirace uvedenými zdroji.
% Kontrola, obsahové porovnání zdrojů a časy: planovani/DOKONCENI_PRIPRAV_13_24.md.
% Učebnice: ucebnice/01_Zakladni_poznatky/Mnoziny/Množiny.tex
% oddíly Rovnost množin, Podmnožiny, Množinové operace a Vennovy diagramy; dle tématu hodiny.
% Realisticky.cz: externi_zdroje/realisticky/01_Zakladni_poznatky/03_Mnoziny/010303_Mnozinove_operace/010303_Lekce.pdf
% Realisticky.cz: externi_zdroje/realisticky/01_Zakladni_poznatky/03_Mnoziny/010304_Vennovy_diagramy/010304_Lekce.pdf
\documentclass[11pt,a4paper]{article}
% Spolecny styl pro ucitelske pripravy MAT (Work/MAT)
% Verze 1.1, 2026-10-23
\NeedsTeXFormat{LaTeX2e}
\ProvidesPackage{priprava}[2026/10/23 v1.1 Sablona ucitelskych priprav MAT]

% Zakladni typografie a cestina
\RequirePackage[T1]{fontenc}
\RequirePackage[utf8]{inputenc}
\RequirePackage[czech]{babel}
\RequirePackage{lmodern}
\RequirePackage{microtype}
\RequirePackage{amsmath,amssymb}
\RequirePackage{xcolor}
\RequirePackage{enumitem}
\RequirePackage[most]{tcolorbox}
\RequirePackage{tabularx,booktabs,array}
\RequirePackage{fancyhdr}
\RequirePackage{needspace}
\RequirePackage{geometry}
\RequirePackage[hidelinks]{hyperref}

\geometry{a4paper,top=17mm,bottom=18mm,left=23mm,right=23mm}
\definecolor{notegray}{gray}{0.39}
\setlength{\parindent}{0pt}
\setlength{\parskip}{2.5pt plus 1pt minus 0.5pt}
\setlist[itemize]{topsep=2pt,itemsep=1pt,parsep=0pt,leftmargin=1.45em}
\setlist[enumerate]{topsep=2pt,itemsep=2pt,parsep=0pt,leftmargin=1.65em}
\setcounter{secnumdepth}{0}
\renewcommand{\arraystretch}{1.12}

% Zahlavi: vlevo rocnik | globalni tema, vpravo cislo: nazev hodiny
% Zapati zustava zcela prazdne.
\pagestyle{fancy}
\fancyhf{}
\renewcommand{\headrulewidth}{0.35pt}
\renewcommand{\footrulewidth}{0pt}
\setlength{\headheight}{14pt}
\newcommand{\nastavPripravu}[4]{%
  \fancyhead[L]{\small #1\enspace\textbar\enspace #2}%
  \fancyhead[R]{\small #3: #4}%
}

% Minimalisticke cernobile boxy. Kazda definice zvlast.
\tcbset{commonlesson/.style={enhanced,breakable=false,colback=white,colframe=black,
  colbacktitle=white,coltitle=black,boxrule=0.45pt,
  arc=0.8mm,outer arc=0.8mm,left=3mm,right=3mm,top=1.6mm,bottom=1.6mm,
  toptitle=1mm,bottomtitle=0.7mm,fonttitle=\bfseries,
  attach boxed title to top left={xshift=2.5mm,yshift=-1.15mm},
  boxed title style={colback=white,colframe=white,boxrule=0pt,left=1mm,right=1mm}}}
\newtcolorbox{zadaniBox}[1][]{commonlesson,title={Zadání},#1}
\newtcolorbox{prikladBox}[1][]{commonlesson,title={Příklad},#1}
\newtcolorbox{definiceBox}[1][]{commonlesson,title={Definice},#1}

% Spolecny box pro zadani a jeho reseni, bez mezery mezi castmi.
% Pouziti:
%   \begin{prikladReseniBox}
%     Zadani ulohy.
%     \tcblower
%     \begin{align*} ... \end{align*}
%   \end{prikladReseniBox}
% Spodni cast ma jemny sedy podklad a delici linku; lze v ni sazet
% vice radku matematiky. Obe casti maji jen jeden spolecny ramecek.
\tcbset{solutionlesson/.style={
  commonlesson,
  enhanced jigsaw,
  breakable,
  colbacklower=black!3,
  segmentation style={solid,draw=black!25,line width=0.35pt},
  fontlower=\small\color{notegray},
  before lower={\noindent\textbf{Řešení:}\enspace},
  leftlower=3mm,rightlower=3mm
}}
\newtcolorbox{prikladReseniBox}[1][]{solutionlesson,title={Příklad},#1}
\newtcolorbox{zadaniReseniBox}[1][]{solutionlesson,title={Zadání},#1}

% Pomucky pro autora pripravy: sedive poznamky; \reseni pro starsi pripravy
\newcommand{\poznamka}[1]{\par\begingroup\small\color{notegray}\itshape
  \textnormal{\textbf{Poznámka pro učitele.}}\enspace #1\par\endgroup}
\newcommand{\reseni}[1]{\par\begingroup\small\color{notegray}
  \textbf{Očekávané řešení:}\enspace #1\par\endgroup}
\newcommand{\blok}[1]{\needspace{5\baselineskip}\subsection*{#1}}
\newcommand{\vyrok}[1]{\uv{#1}}

\usepackage{tikz}
\usetikzlibrary{patterns,calc,angles,quotes}
\nastavPripravu{1. ročník}{Množiny}{15}{Sjednocení a průnik}
\begin{document}
\blok{Dvě podmínky: současně, nebo alespoň jedna}
V evidenci součástek $A$ označuje díly z hliníku a $B$ díly s otvorem.
Výběr „z hliníku a s otvorem“ vyžaduje obě podmínky.
Výběr „z hliníku nebo s otvorem“ zahrnuje i díly splňující obě podmínky.
Navazujeme na spojky $\land$ a $\lor$ z logiky a na oblasti diagramu.

\begin{definiceBox}
\textbf{Průnik} $A\cap B$ tvoří právě společné prvky:
$x\in A\cap B\Leftrightarrow (x\in A\land x\in B)$.
\end{definiceBox}
\begin{definiceBox}
\textbf{Sjednocení} $A\cup B$ tvoří prvky alespoň jedné z obou množin:
$x\in A\cup B\Leftrightarrow (x\in A\lor x\in B)$.
\end{definiceBox}
\begin{center}
\begin{tikzpicture}[scale=.75]
\begin{scope}
\clip (-.8,0) circle(1.3); \clip (.8,0) circle(1.3);
\fill[pattern=north east lines] (-2.2,-1.5) rectangle (2.2,1.5);
\end{scope}
\draw (-.8,0) circle(1.3); \draw (.8,0) circle(1.3);
\node at (-1.3,1.5) {$A$}; \node at (1.3,1.5) {$B$};
\node at (0,-1.75) {$A\cap B$: obě podmínky};
\begin{scope}[xshift=6.5cm]
\fill[black!15] (-.8,0) circle(1.3); \fill[black!15] (.8,0) circle(1.3);
\draw (-.8,0) circle(1.3); \draw (.8,0) circle(1.3);
\node at (-1.3,1.5) {$A$}; \node at (1.3,1.5) {$B$};
\node at (0,-1.75) {$A\cup B$: alespoň jedna};
\end{scope}
\end{tikzpicture}
\end{center}
\begin{prikladReseniBox}[title={P1: výčet i kontrola diagramem},breakable=false]
Pro $A=\{-2,0,2,4\}$ a $B=\{0,1,4,7\}$ urči $A\cap B$ a $A\cup B$.
Která z hledaných množin obsahuje $0$?
\tcblower
$A\cap B=\{0,4\}$, $A\cup B=\{-2,0,1,2,4,7\}$.
Nula patří do obou výsledků. Ve sjednocení společné prvky neopakujeme.
V diagramu by jen v $A$ ležely $-2,2$, ve společné části $0,4$ a jen v $B$ $1,7$.
\end{prikladReseniBox}

\begin{definiceBox}
Množiny jsou \textbf{disjunktní}, jestliže nemají společný prvek:
$A\cap B=\varnothing$.
\end{definiceBox}
\begin{prikladReseniBox}[title={P2: číselný obor je podstatný},breakable=false]
V $\mathbb Z$ označ $S$ sudá a $L$ lichá čísla.
Urči $S\cap L$, $S\cup L$ a $\mathbb N\cap\mathbb Z$.
\tcblower
$S\cap L=\varnothing$: celé číslo nemůže být sudé i liché.
$S\cup L=\mathbb Z$, protože každé celé číslo patří právě do jedné skupiny.
$\mathbb N\cap\mathbb Z=\mathbb N$, neboť $\mathbb N\subseteq\mathbb Z$.
Nula patří do $S$, ale ne do naší $\mathbb N$.
\end{prikladReseniBox}

\blok{Co plyne přímo z významu operací}
Platí $A\cap B\subseteq A\subseteq A\cup B$ a stejně pro $B$.
Přidání požadavku může výběr zúžit, možnost volby jej může rozšířit;
nemusí však jít o vlastní zahrnutí. Je-li $A\subseteq B$, pak
$A\cap B=A$ a $A\cup B=B$. Obrácené pořadí v průniku i sjednocení dává stejný výsledek.

\begin{prikladReseniBox}[title={P3: tři množiny a závorky},breakable=false]
Nechť $A=\{1,2,3\}$, $B=\{3,4\}$, $C=\{2,4,5\}$.
Porovnej $(A\cup B)\cap C$ a $A\cup(B\cap C)$.
\tcblower
Nejdříve $A\cup B=\{1,2,3,4\}$, tedy $(A\cup B)\cap C=\{2,4\}$.
Dále $B\cap C=\{4\}$, takže $A\cup(B\cap C)=\{1,2,3,4\}$.
Výsledky se liší například v prvku $1$. Závorky nelze svévolně přemístit.
\end{prikladReseniBox}
\poznamka{Jádro: definice, P1–P3 a S1 a)–c), přibližně 38 min práce;
organizace, shrnutí a domácí zadání dohromady 7 min. S1 d) lze dokončit doma.
K výpočtu vždy nejprve přečti zadání slovy; spodní části ukazuj až po pokusu žáků.}

\blok{Samostatné procvičení}
\begin{zadaniReseniBox}[title={S1: od výčtu ke vztahu}]
\begin{enumerate}[label=\alph*)]
\item Pro $A=\{-3,-1,2\}$, $B=\{-1,0,2,5\}$ vypiš průnik a sjednocení.
\item Urči $A\cap\varnothing$, $A\cup\varnothing$, $A\cap A$, $A\cup A$ a vysvětli je slovně.
\item Jsou $\{x\in\mathbb Z\mid x<0\}$ a $\{x\in\mathbb Z\mid x\ge0\}$
disjunktní? Co je jejich sjednocení?
\item Najdi dvě neprázdné různé množiny, pro které $A\cap B=A$.
Zdůvodni obecnou podmínku.
\end{enumerate}
\tcblower
a) $\{-1,2\}$ a $\{-3,-1,0,2,5\}$.
b) $\varnothing,A,A,A$: žádný společný prvek s prázdnou množinou,
žádný nový při jejím připojení, žádná změna při porovnání se sebou.
c) Ano, průnik je prázdný a sjednocení $\mathbb Z$; nula patří do druhé množiny.
d) Např. $A=\{1\}$ a $B=\{1,2\}$. Obecně právě $A\subseteq B$:
každý prvek $A$ musí být také v $B$, aby jej průnik zachoval.
\end{zadaniReseniBox}

\blok{Volitelná rezerva}
\begin{zadaniReseniBox}[title={R1: obě podmínky nemusí znamenat obě množiny stejné}]
Pro $U=\{1,\ldots,9\}$, $A=\{1,3,5,7,9\}$, $B=\{3,6,9\}$ a $C=\{1,2,3\}$
urči $A\cap B\cap C$, $(A\cup C)\cap B$ a $A\cup B\cup C$.
Potom rozhodni, zda z $A\cap B=\varnothing$ plyne $A\cup B=\varnothing$.
\tcblower
Postupně $\{3\}$, $\{3,9\}$, $\{1,2,3,5,6,7,9\}$.
Poslední obecné tvrzení neplatí: $A=\{1\}$ a $B=\{2\}$
mají prázdný průnik, sjednocení $\{1,2\}$ však prázdné není.
\end{zadaniReseniBox}
\poznamka{Rezerva 8–10 minut. Při kontrole kombinovaných výrazů požaduj
aspoň jeden mezivýsledek; nedoplňuj operace podle tvaru značky.}

\blok{Domácí procvičení}
\begin{zadaniReseniBox}[title={D1: 10–15 minut}]
\begin{enumerate}[label=\alph*)]
\item Pro $K=\{0,2,4,6\}$ a $L=\{1,2,3,4\}$ urči $K\cap L$, $K\cup L$.
\item Urči $\mathbb N\cup\mathbb Z$ a $\{0\}\cap\mathbb N$.
\item Pro $A=\{a,b\}$, $B=\{b,c\}$, $C=\{c,d\}$ urči
$(A\cup B)\cap C$ a $A\cap(B\cup C)$.
\item Najdi neprázdné disjunktní podmnožiny $U=\{1,2,3,4\}$,
jejichž sjednocení je celé $U$.
\item Žák tvrdí: sjednocení má vždy více prvků než každá původní množina.
Rozhodni a uveď příklad.
\end{enumerate}
\tcblower
a) $\{2,4\}$ a $\{0,1,2,3,4,6\}$.
b) $\mathbb Z$ a $\varnothing$.
c) $\{c\}$ a $\{b\}$.
d) Například $\{1,3\}$ a $\{2,4\}$; nemají společný prvek.
e) Neplatí, např. pro $A=B=\{1\}$ zůstane sjednocení $\{1\}$.
Obecně zaručujeme zahrnutí, nikoli větší počet.
\end{zadaniReseniBox}
\end{document}
